\begin{thebibliography}{99}

% ---------- Sources institutionnelles et marché ----------

\bibitem{france_assureurs_2025}
FRANCE ASSUREURS. \textit{Présentation --- Conférence de presse 2026}. (Mars 2026). Consulté sur \url{https://www.franceassureurs.fr/wp-content/uploads/20260325_presentation-conference-de-presse-2.pdf}

\bibitem{france_assureurs_donnees_cles}
FRANCE ASSUREURS. \textit{Données clés 2024}. (2024). Consulté sur \url{https://www.franceassureurs.fr/wp-content/uploads/donnees-cles-2024.pdf}

\bibitem{insurance_europe_fraud}
INSURANCE EUROPE. \textit{The impact of insurance fraud}. Consulté sur \url{https://www.insuranceeurope.eu/mediaitem/0bf0af82-e7ef-4439-a763-d7862859d421/The+impact+of+insurance+fraud.pdf}

\bibitem{alfa_news_48}
ALFA (Agence de Lutte contre la Fraude à l'Assurance). \textit{ALFA News n\textsuperscript{o}~48}. (Septembre 2024). Consulté sur \url{https://www.alfa.asso.fr/alfa-news-48-septembre-2024/}

\bibitem{alfa_news_39}
ALFA (Agence de Lutte contre la Fraude à l'Assurance). \textit{ALFA News n\textsuperscript{o}~39}. (Juin 2023). Consulté sur \url{https://www.alfa.asso.fr/alfa-news-39-juin-2023-2/}

\bibitem{argus_cout_2_5md}
L'ARGUS DE L'ASSURANCE. \textit{Fraude à l'assurance : une facture de 2,5~Md\euro{} en dommages}. Consulté sur \url{https://www.argusdelassurance.com/institutions/fraude-a-l-assurance-une-facture-de-2-5-md-en-dommages-en-2014-alfa.96060}

\bibitem{argus_cout_50}
L'ARGUS DE L'ASSURANCE. \textit{Fraude à l'assurance : un coût de 50~\euro{} pour les assurés}. Consulté sur \url{https://www.argusdelassurance.com/acteurs/fraude-a-l-assurance-un-cout-de-50-pour-les-assures.60931}

\bibitem{argus_fraude_2025}
L'ARGUS DE L'ASSURANCE. \textit{Près d'un milliard d'euros de fraudes à l'assurance détecté en 2025}. Consulté sur \url{https://www.argusdelassurance.com/assurance-de-personnes/des-filieres-de-larnaque-utilisent-lintelligence-artificielle-pres-dun-milliard-deuros-de-fraudes-a-lassurance-detecte-en-2025.VRVQHSDANZGGZL5GXH7FHW3JU4.html}

\bibitem{itesoft_fraude}
ITESOFT. \textit{Fraude à l'assurance : les chiffres clés (ALFA)}. Consulté sur \url{https://www.itesoft.com/fr/blog/fraude-assurance-chiffres-cles-alfa/}

% ---------- Cadre réglementaire ----------

\bibitem{acpr_gouvernance_ia}
ACPR (Autorité de Contrôle Prudentiel et de Résolution). \textit{Gouvernance des algorithmes d'intelligence artificielle dans le secteur financier}. (Juin 2020). Consulté sur \url{https://acpr.banque-france.fr/system/files/2025-02/20200612_gouvernance_evaluation_ia.pdf}

\bibitem{rgpd_art22}
UNION EUROPÉENNE. \textit{Règlement général sur la protection des données (RGPD), article 22 --- Décision individuelle automatisée}. Consulté sur \url{https://gdpr-text.com/fr/read/article-22/}

\bibitem{ai_act}
UNION EUROPÉENNE. \textit{Règlement (UE) 2024/1689 établissant des règles harmonisées concernant l'intelligence artificielle (AI Act)}. Journal officiel de l'Union européenne, 2024.

\bibitem{code_assurances}
CODE DES ASSURANCES. \textit{Articles L.121-1, L.113-2, L.113-8 et L.113-9}. Legifrance.

\bibitem{code_penal_313}
CODE PÉNAL. \textit{Article 313-1 --- Escroquerie}. Legifrance.

% ---------- Fraude au système d'immatriculation (SIV) ----------

\bibitem{siv_assemblee}
ASSEMBLÉE NATIONALE. \textit{Question écrite n\textsuperscript{o}~4793 --- Fraude au système d'immatriculation des véhicules}. Consulté sur \url{https://questions.assemblee-nationale.fr/q17/17-4793QE.htm}

\bibitem{siv_senat}
SÉNAT. \textit{Question écrite n\textsuperscript{o}~02893 --- Fraude au système d'immatriculation des véhicules}. (2025). Consulté sur \url{https://www.senat.fr/questions/base/2025/qSEQ250102893.html}

% ---------- Données et outils ----------

\bibitem{kaggle_dataset}
AHLUWALIA Saksham. \textit{Car Insurance Fraud Detection Dataset}. Kaggle. Consulté sur \url{https://www.kaggle.com/datasets/ahluwaliasaksham/car-insurance-fraud-detection-dataset}

\bibitem{anthropic_claude}
ANTHROPIC. \textit{Claude --- Assistant conversationnel d'intelligence artificielle}. Consulté sur \url{https://www.anthropic.com}

% ---------- Références méthodologiques (machine learning) ----------

\bibitem{breiman_rf}
BREIMAN Leo. « Random Forests ». \textit{Machine Learning}, volume 45, n\textsuperscript{o}~1, 2001, p.~5-32.

\bibitem{friedman_gbm}
FRIEDMAN Jerome H. « Greedy Function Approximation: A Gradient Boosting Machine ». \textit{Annals of Statistics}, volume 29, n\textsuperscript{o}~5, 2001, p.~1189-1232.

\bibitem{chen_xgboost}
CHEN Tianqi, GUESTRIN Carlos. « XGBoost: A Scalable Tree Boosting System ». \textit{Proceedings of the 22nd ACM SIGKDD International Conference on Knowledge Discovery and Data Mining}, 2016, p.~785-794.

\bibitem{liu_isolation_forest}
LIU Fei Tony, TING Kai Ming, ZHOU Zhi-Hua. « Isolation Forest ». \textit{Proceedings of the 8th IEEE International Conference on Data Mining}, 2008, p.~413-422.

\bibitem{lundberg_shap}
LUNDBERG Scott M., LEE Su-In. « A Unified Approach to Interpreting Model Predictions ». \textit{Advances in Neural Information Processing Systems (NeurIPS)}, volume 30, 2017.

\bibitem{blondel_louvain}
BLONDEL Vincent D., GUILLAUME Jean-Loup, LAMBIOTTE Renaud, LEFEBVRE Etienne. « Fast unfolding of communities in large networks ». \textit{Journal of Statistical Mechanics: Theory and Experiment}, 2008.

\end{thebibliography}